% Master quant CV. Layout: templates/cv.cls. Section order follows tracks/quant.md.
% Replace the example content with your own, or ask Claude to build it from your profile.
% End every bullet with "% src: <entry-id>#F<n>".
\documentclass{cv}

\cvname{{{NAME}}}
\cvlocation{{{CITY}}}
\cvemail{{{EMAIL}}}
\cvphone{{{PHONE}}}{{{PHONE_LINK}}}
\cvlinkedin{{{LINKEDIN_URL}}}{{{NAME}}}

\begin{document}

\makecvheader

\section{Key Highlights}
\begin{bullets}
  \item <Your strongest quantitative project, one line.> % src: example#F1
\end{bullets}

\section{Education}
\education{<University>}{<Degree> (<year> of <n> years)}{Aug 2024 -- Present}
\begin{bullets}
  \item \textbf{Average grade:} <grade> % src: example#F2
  \item \textbf{Coursework:} Statistics, <probability / optimisation / ML courses>. % src: example#F3
\end{bullets}

\section{Experience}
\entry{<Job title>}{<Organisation>}{<City, Country>}{June 2026 -- Aug 2026}
\begin{bullets}
  \item <Quantitative or analytical work, with a concrete detail.> % src: example#F4
\end{bullets}

\section{Projects}
\entry{<Modelling or trading project>}{Personal project}{}{2025 -- Present}
\begin{bullets}
  \item <Data and models used.> % src: example#F5
  \item <How you validated it (out-of-sample, bias controls). No best-of-N results.> % src: example#F6
\end{bullets}

\section{Awards}
\begin{bullets}
  \item <Competition, placement, year.> % src: example#F7
\end{bullets}

\section{Skills \& Other}
\begin{skilltable}
  \row{Programming}{Python (pandas, NumPy), C++, SQL (basic)}
  \row{ML \& data}{PyTorch, scikit-learn}
  \row{Languages}{<language> (native), English (fluent)}
  \row{Interests}{<interests>}
\end{skilltable}

\end{document}
